% !PRINTED 220
\weberpage{220}
%<<CONT

\weq{5}{y = t_1\Bigl(a_0x + \frac{a_1}{3}\Bigr) + t_0\Bigl(a_0x^2 + a_1x + \frac{2a_2}{3}\Bigr)}

results for $y$. If we write the equation in the form

\[ y^3 + P_2y + P_3 = 0, \]

then $P_2, P_3$ are covariants of $f$, and indeed $P_2$ of the second, $P_3$ of the third order, both in $t$ and in $a$.

Now since $P_2$ and $P_3$ are covariants of the cubic form, by \S.\,62 and 63 they can differ from the two functions defined there,

\[ H(t_1,t_0), \qquad Q(t_1,t_0) \]

only by constant, that is, numerical factors. These constant factors can be determined by some special assumption.

We may, for example, assume

\[ t_1 = 1, \quad t_0 = 0, \quad a_0 = 1, \quad a_1 = 0, \]

and then $y=x$, hence

\[ P_2 = a_2, \qquad P_3 = a_3. \]

On the other hand, by the formulas (2), (3), (7) given in \S.\,62, this particular assumption yields

\[ H(1,0) = 3a_2, \qquad Q(1,0) = 27a_3, \]

from which one concludes in general

\[ P_2 = \frac{1}{3}H(t_1,t_0), \qquad P_3 = \frac{1}{27}Q(t_1,t_0), \]

so that we obtain for $y$ the cubic equation

\weq{6}{y^3 + \frac{1}{3}H(t_1,t_0)y + \frac{1}{27}Q(t_1,t_0) = 0.}

The discriminant $\Delta$ of this cubic equation is

\[ \Delta = -\frac{1}{27}(4H^3 + Q^2), \]

hence, applying the relation \S.\,62, (12),

\weq{7}{4H^3 + Q^2 + 27Df^2 = 0}

\weq{8}{\Delta = Df(t_1,t_0)^2,}

if $D$ is the discriminant of the given cubic equation, in agreement with the general results of the preceding section.

If we wish to base the solution of the cubic equation on this, we must first, following the prescriptions of the preceding section, express $x$ rationally in terms of $y$. We set

\weq{9}{a_0f(t_1,t_0)x = Q_0 + Q_1y + Q_2y^2.}
